\chapter{Black Hole Thermodynamics}
\label{ch:bh-thermodynamics}
Before developing the full quantum field theory in curved spacetime,
we establish that black holes behave as thermodynamic objects even at
the semiclassical level. This was historically the first hint that
black holes and thermodynamics are deeply connected, and it sets up the
physical question that the Hawking calculation answers: if a black hole
has entropy proportional to its horizon area, what is the temperature
associated with that entropy, and is there a corresponding thermal
radiation? The gedanken experiments of this chapter fix the
Bekenstein-Hawking entropy $S_{BH} = A/4$ and the generalized second
law, providing the thermodynamic framework within which Hawking
radiation will later be understood.

The idea that black holes carry entropy and behave as thermodynamic
objects was proposed by Bekenstein on the basis of a simple but
compelling analogy. We develop this analogy here and subject it to
quantitative tests.

\section{The Bekenstein Entropy Conjecture}

Consider throwing a box of gas with entropy $S_p$ into a black hole.
The gas disappears behind the horizon and its entropy is lost to an
outside observer. But the second law of thermodynamics forbids the
total entropy of the universe from decreasing. This suggests that the
black hole itself must carry entropy to compensate.

Two facts from general relativity support this. First, vacuum solutions
to the Einstein equations with regular horizons, that is, black holes,
are determined by just two parameters: mass $M$ and angular momentum
$J$ (plus charge $Q$ for non-vacuum solutions). This is the
no-hair theorem. Second, the area of a black hole horizon never
decreases; it only grows. For a Schwarzschild black hole,
\begin{equation}
	A = 4\pi r_h^2 = 4\pi(2M)^2 = 16\pi M^2.
\end{equation}
The area has the same monotonicity property we expect of entropy.
Bekenstein therefore proposed that the black hole entropy is
proportional to the horizon area,
\begin{equation}
	S_{BH} = \alpha\cdot A_{\text{horizon}}
	\overset{\text{Schwarzschild}}{=} 16\pi\alpha M^2,
\end{equation}
where $\alpha$ is a coefficient to be determined. The area $A$ here is
the proper area, coming from the angular part of the metric
$r^2(\tilde u,\tilde v)\,d\Omega^2$, and is therefore
coordinate-independent.

\section{The Qubit Gedanken Experiment}

To test the conjecture and fix $\alpha$, consider lowering a qubit
inside a box into a Schwarzschild black hole using a string. The qubit
has entropy $S_p = \ln 2$ (it can be spin up or spin down in a mixed
state) and mass $\mu$. The question is whether the total entropy
$S_{\text{tot}} = S_{BH} + S_{\text{outside}}$ can decrease in this
process.

% [FIGURE: qubit being lowered into black hole on a string]

The qubit has four-momentum $p_\mu$ with
\begin{equation}
	\mu^2 = -g^{\mu\nu}p_\mu p_\nu.
\end{equation}
Since we lower the qubit radially, $p_r = 0$, and the energy measured
at infinity is
\begin{equation}
	E = \mu\sqrt{1-\frac{2M}{r}},
\end{equation}
which is simply gravitational redshift. By energy conservation, the
change in black hole mass when the qubit falls in at radius $r$ is
$\delta M = E$. It appears that by bringing the qubit arbitrarily
close to the horizon we can make $\delta M$, and hence $\delta S_{BH}$,
arbitrarily small, in apparent contradiction with the second law.

The resolution comes from the uncertainty principle. To control the
number of particles in the box we need the uncertainty in energy to
satisfy $\delta E < \mu$, which means the particle can only be
localized to a region larger than its Compton wavelength,
\begin{equation}
	\delta l \geq \frac{1}{\mu}.
\end{equation}
The proper distance $\delta l$ and the coordinate distance $\delta r$
near the horizon are related by
\begin{equation}
	dl^2 = \frac{dr^2}{1-\frac{2M}{r}}
	\implies l = \int_{r_h}^{r_h+\delta r}
	\frac{dr}{\sqrt{1-\frac{2M}{r}}}
	\simeq 2\sqrt{2M\delta r},
\end{equation}
so the constraint $\delta l \geq 1/\mu$ gives
$\delta r \geq 1/(8M\mu^2)$. Near the horizon at
$r = r_h + \delta r$,
\begin{equation}
	E^2 = \mu^2\left(1-\frac{2M}{r_h+\delta r}\right)
	\simeq \mu^2\frac{\delta r}{2M},
\end{equation}
and so $E = \mu\sqrt{\delta r/2M}$. Using the minimum $\delta r$,
\begin{equation}
	E \geq \mu\cdot\frac{1}{\sqrt{2M}}\cdot
	\frac{1}{\sqrt{8M\mu^2}} = \frac{1}{4M},
\end{equation}
giving the fundamental result
\begin{equation}
	\boxed{\delta M \geq \frac{1}{4M}.}
\end{equation}
The change in horizon area is then
\begin{equation}
	\delta A_h = \delta(16\pi M^2)
	= 32\pi M\,\delta M \geq 8\pi > 0,
\end{equation}
in Planck units. This is finite, not zero, so the entropy of the black
hole does increase. Dropping one qubit of entropy $\ln 2$ corresponds
to a minimum area increase $\delta A_h \geq 8\pi$, so the minimum
entropy increase satisfies
\begin{equation}
	\frac{\delta S_{BH}}{\delta A_h} = \alpha
	\implies \alpha = \frac{\ln 2}{8\pi}.
\end{equation}
The correct coefficient, fixed by Hawking's calculation, is
$\alpha = 1/4$, giving the Bekenstein-Hawking entropy
\begin{equation}
	\boxed{S_{BH} = \frac{A_h}{4}.}
\end{equation}
Our estimate gives the right order of magnitude and confirms the
positivity of $\delta S_{\text{tot}}$.

\section{The Generalized Second Law}

We now verify the generalized second law $\delta S_{\text{tot}} \geq 0$
for two further scenarios.

\subsection*{Geodesic Infall with Angular Momentum}

Consider the same setup but now the qubit falls along a geodesic with
angular momentum $\tilde L$. For a massive particle in the Schwarzschild
geometry, the geodesic equation gives
\begin{equation}
	\dot r^2 = \tilde E^2
	- \left(1-\frac{2M}{r}\right)
	\left(1+\frac{\tilde L^2}{r^2}\right)
	\equiv \tilde E^2 - V_{\text{eff}}(r),
\end{equation}
where $\tilde E = E/\mu$ and $\tilde L = L/\mu$ are the specific energy
and angular momentum. The energy deposited into the black hole is
\begin{equation}
	E = \mu\left[\dot r^2
		+ f(r)\left(1+\frac{\tilde L^2}{r^2}\right)\right]^{1/2},
\end{equation}
and the changes in mass and angular momentum are $\delta M = E$,
$\delta J = \tilde L$. The change in black hole entropy is
\begin{equation}
	\delta S_{BH} = 32\pi\alpha M^2\mu\cdot E
	= 32\pi\alpha M^2\mu^2
	\left[\dot r^2 + f(r)\left(1+\frac{\tilde L^2}{r^2}\right)
		\right]^{1/2}.
\end{equation}
Near the horizon, $f(r_h+\delta r) = \delta r/2M$ and
$\dot r^2\sim 1/(8M\mu^2\delta r)$, so
\begin{equation}
	E\sim\frac{1}{\sqrt{8M\delta r}}
	\geq \frac{1}{\sqrt{8M\cdot\frac{1}{8M\mu^2}}}
	= \mu,
\end{equation}
giving $\delta S_{BH} \geq 32\pi\alpha M^2\mu^2$. With $\alpha = 1/4$
this gives $\delta S_{BH} \geq 8\pi M^2\mu^2 \gg \ln 2$, so
\begin{equation}
	\delta S_{\text{tot}}
	= \delta S_{BH} - \ln 2
	\geq 8\pi M^2\mu^2 - \ln 2 > 0,
\end{equation}
regardless of the angular momentum of the infalling particle.

\subsection*{Box of Radiation}

Consider a black hole and a box of radiation at temperature $T$ with
energy $E_{\text{box}}$ and entropy $S_{\text{box}} = E_{\text{box}}/T$.
A machine lowers the box to near the horizon, opens it, lets the
radiation fall in, then raises the empty box back. The energy deposited
into the black hole at $r = r_h + \delta r$ is
\begin{equation}
	\delta M = E_{\text{box}}\sqrt{f(r_h+\delta r)}
	\simeq E_{\text{box}}\sqrt{\frac{\delta r}{2M}}.
\end{equation}
The change in black hole entropy is
\begin{equation}
	\delta S_{BH} = 32\pi\alpha M^2\,\delta M
	= 32\pi\alpha M^2 E_{\text{box}}\sqrt{\frac{\delta r}{2M}},
\end{equation}
and the entropy lost from outside is $\delta S_{\text{outside}}
	= -S_{\text{box}} = -E_{\text{box}}/T$. The box must be at least as
large as the thermal wavelength of the radiation, $\delta r \geq 1/T$,
so
\begin{equation}
	\sqrt{\frac{\delta r}{2M}} \geq \frac{1}{\sqrt{2MT}},
\end{equation}
and therefore
\begin{equation}
	\delta S_{BH} \geq \frac{32\pi\alpha M^2 E_{\text{box}}}{\sqrt{2MT}}.
\end{equation}
With $\alpha = 1/4$ and $T_{BH} = 1/(8\pi M)$,
\begin{equation}
	\delta S_{BH} \geq \frac{8\pi M^2 E_{\text{box}}}{\sqrt{2MT}}
	\gg \frac{E_{\text{box}}}{T} = S_{\text{box}},
\end{equation}
so $\delta S_{\text{tot}} = \delta S_{BH} - S_{\text{box}} \geq 0$.
The machine cannot violate the second law.

\subsection*{Charged Black Hole}

For a black hole with charge $Q$, the metric is the
Reissner-Nordstr\"{o}m metric,
\begin{equation}
	ds^2 = -f(r)\,dt^2 + \frac{dr^2}{f(r)} + r^2\,d\Omega^2,
	\qquad f(r) = 1 - \frac{2M}{r} + \frac{Q^2}{r^2},
\end{equation}
with electric potential $A_0 = Q/r$. The horizon is at
\begin{equation}
	r_h = M + \sqrt{M^2 - Q^2},
\end{equation}
and the horizon area is
\begin{equation}
	A_h = 4\pi r_h^2
	= 4\pi\left(M+\sqrt{M^2-Q^2}\right)^2.
\end{equation}
Now consider lowering a charged particle with mass $\mu$ and charge $q$
to near the horizon. The conserved energy for a charged particle in
this background is $E = \mu\sqrt{f(r)+\dot r^2} + qA_0$, so the
energy and charge deposited into the black hole are
$\delta M = E_\infty$ and $\delta Q = q$. The change in horizon area
follows from differentiating $f(r_h) = 0$,
\begin{equation}
	\frac{\partial r_h}{\partial M}
	= \frac{r_h^2}{r_h^2-Q^2}, \qquad
	\frac{\partial r_h}{\partial Q}
	= \frac{-Qr_h}{r_h^2-Q^2},
\end{equation}
giving
\begin{equation}
	\delta A_h = \frac{8\pi r_h^2}{r_h^2-Q^2}
	\left(\delta M - \Phi_h\,\delta Q\right),
\end{equation}
where $\Phi_h = Q/r_h$ is the horizon electric potential. The condition
$\delta M - \Phi_h\delta Q > 0$ ensures the particle actually falls in
and is not repelled. Using the uncertainty principle as before,
$\delta r \geq 1/(2\mu)$, and combining with the entropy lost outside,
\begin{equation}
	\delta S_{\text{tot}} = \delta S_{BH} - S_{\text{box}}
	= \frac{2\pi r_h^2}{r_h^2-Q^2}
	\left(\delta M - \Phi_h\delta Q\right) - \frac{E_{\text{box}}}{T}
	\geq 0,
\end{equation}
so the generalized second law holds for the Reissner-Nordstr\"{o}m
black hole as well.
